%%%%%%%%%%%%%%%%%%%%%%%%%%%%%%%%%%%
% Academia.edu Submission Template
% Body content duplicated from latex/master.tex into content.tex
% Original RSC format in latex/master.tex is unchanged
%%%%%%%%%%%%%%%%%%%%%%%%%%%%%%%%%%%

\documentclass[12pt]{Definitions/academia}

\usepackage{hyperref}
\usepackage[numbers,sort&compress]{natbib}
\usepackage[nobreak=false]{mdframed}
\usepackage{mathtools}
\usepackage{enumitem}
\usepackage{float}

% Commands used in RSC format that need to be defined here
\providecommand{\balance}{}

\allowdisplaybreaks

\begin{document}

\articletype{Research Article}
\proofstage{Preprint}
\jName{Physics}

\title{Unifying General and Special Relativity: Demonstrating a Unified Inertial Origin of Relativistic Time Dilation}

\author{$^{1*}$Thomas Damon DeGerlia}

\maketitle

\affiliation{$^{1}$DeGerlia Expert Consulting, 3000 Lawrence Street, Denver, CO, United States of America}

\noindent $^{*}$Correspondence: Thomas Damon DeGerlia
\email{$^{*}$tom.degerlia@tomdegerlia.com}

\section*{Abstract}
Inertial Relativity and its key formula the Space-Time Equivalence (STE)\citep{degerlia202501} hypothesize that there is a universal mathematical relationship between time and the spatial distribution of matter: \textit{The characteristic length scale factor that determines time dilation between any two physical systems about any respective axes is universally equal to the fifth root of the inverse ratio of their moments of inertia about the same axes.} A central prediction of this hypothesis is that General and Special Relativity are edge cases of the STE. In our prior paper \citep{degerlia202505}, we began validating this prediction by arriving at General Relativity calculations with exact precision utilizing only principles of relativistic inertia. This paper completes the validation by demonstrating the mathematical equivalence between the STE, General Relativity, and Special Relativity. In the context of moment of inertia about an axis, we demonstrate that the STE is mathematically equivalent to Schwarzschild gravitational time dilation. In the context of linear inertia (mass), the STE states an inverse third root relationship with inertia (mass), which we demonstrate is mathematically equivalent to Lorentz time dilation. Both results confirm the original prediction: Special and General Relativity emerge as edge cases of the same underlying phenomenon, Inertial Relativity: Special Relativity at the extremes of velocity and General Relativity at the extremes of mass, unifying the two under a single universal relativistic framework.

\keywords{modern physics; inertial relativity; space-time equivalence; time dilation; general relativity; special relativity; moment of inertia; unification}

\input{content}

\end{document}
